% 本文件由 03_代码/p4_tables.py 自动生成，请勿手工修改。
% 数据源：06_支撑材料/p4_results.json

\begin{table}[htbp]
  \centering
  \caption{附件 4 实时电价的结构分解（$P=F+R$，$F$ 为附件 1 的固定曲线，$R$ 为残差）}
  \label{tab:p4-struct}
  \footnotesize
  \begin{tabular}{lrr}
    \toprule
    项目 & 数值 & 说明 \\
    \midrule
    逐年逐段均值 $\overline{P}$ 与 $F$ 的最大偏差 & $5.15e-05$ & 两者\textbf{骨架与水平相同} \\
    $F$ 的年平均 & 0.7662 & 元/kWh \\
    $\overline{P}$ 的年平均 & 0.7662 & 元/kWh，与 $F$ 相差不足 $10^{-6}$ \\
    $R$ 与 $F$ 的相关系数 & $-0.0000$ & 正交，故 $R$ 就是``波动''本身 \\
    \midrule
    $R$ 的周一均值 & +0.0529 & 元/kWh \\
    $R$ 的周二均值 & +0.0549 & 元/kWh \\
    $R$ 的周三均值 & +0.0576 & 元/kWh \\
    $R$ 的周四均值 & +0.0542 & 元/kWh \\
    $R$ 的周五均值 & -0.1372 & 元/kWh \\
    $R$ 的周六均值 & -0.1364 & 元/kWh \\
    $R$ 的周日均值 & +0.0530 & 元/kWh \\
    \midrule
    周五、周六平均 & -0.1368 & 元/kWh，\textbf{偏低} \\
    其余五天平均 & +0.0545 & 元/kWh \\
    两者之差 & +0.1913 & 元/kWh，相当于均价的 25.0\% \\
    星期结构解释的 $R$ 方差 & 41.7\% & 其余为无规律的日内噪声 \\
    ``同星期加权''恒等式最大残差 & $0.00$ & 见式~\eqref{eq:p4-price} \\
    \bottomrule
  \end{tabular}
  \par\vspace{2pt}
  \begin{minipage}{.92\textwidth}\footnotesize
    注：附件 4 的电价并非与附件 1 无关的另一条曲线，而是\textbf{在附件 1的 144 段骨架上叠加了一个以星期为周期、均值为零的偏移}。这一条决定了后文对照实验的读法：两者年平均值相同，因此任何费用差异都来自\textbf{波动本身}，不掺杂价格水平的变化。周五、周六电价系统性偏低，是残差里唯一可利用的规律，也正是同星期加权预测器的依据。
  \end{minipage}
\end{table}

\begin{table}[htbp]
  \centering
  \caption{价格预测的四个发布版本在各自动态时刻更新后的精度}（框架 6.3 节）
  \label{tab:p4-skill}
  \footnotesize
  \begin{tabular}{lrrrr}
    \toprule
    发布版本 & 覆盖时段 & 平均绝对误差 & 相对均价 & 共同时段误差 \\
    \midrule
    0点版本 & 0--144 & 0.0441 & 5.82\% & 0.0419 \\
    6点版本 & 36--144 & 0.0458 & 6.04\% & \textbf{0.0386} \\
    12点版本 & 72--144 & 0.0433 & 5.71\% & 0.0396 \\
    18点版本 & 108--144 & 0.0398 & 5.25\% & 0.0398 \\
    \midrule
    对照：全天用同一个均价 & 全天 & 0.2774 & 36.62\% & --- \\
    \bottomrule
  \end{tabular}
  \par\vspace{2pt}
  \begin{minipage}{.92\textwidth}\footnotesize
    注：``覆盖时段''是时段索引范围，第 $h$ 个版本只能修正它发布之后的时段，因此四个版本的``平均绝对误差''覆盖的时段\textbf{互不相同}（$0{:}00$ 版覆盖全天 $144$ 段，$18{:}00$ 版只覆盖最后 $36$ 段），\textbf{直接比大小是错的}。最后一列把四个版本都限制在 $18{:}00$--$24{:}00$（四个版本都覆盖该段）上求平均绝对误差，才是可比的版本间排序；加粗者为其中的最优。最后一行的``全天用同一个均价''是没有任何日内形状信息的预测，它的误差远大于四个版本，说明形状信息确有价值。
  \end{minipage}
\end{table}

\begin{table}[htbp]
  \centering
  \caption{价格预测候选方案在 2025 年 2 月 1 日--12 月 31 日的平均绝对误差}
  \label{tab:p4-forecast}
  \footnotesize
  \begin{tabular}{lrr}
    \toprule
    预测方案 & 平均绝对误差/(元/kWh) & 相对均价 \\
    \midrule
    \textbf{同星期加权 0.6/0.3/0.1} & 0.0441 & 5.82\% \\
    同星期加权 0.5/0.3/0.2 & 0.0446 & 5.88\% \\
    同星期等权 1/3 & 0.0466 & 6.15\% \\
    同星期(上周同一天) & 0.0472 & 6.22\% \\
    近 7 天滚动均值 & 0.0827 & 10.91\% \\
    近 14 天滚动均值 & 0.0838 & 11.07\% \\
    昨天同时段 & 0.0843 & 11.13\% \\
    近 28 天滚动均值 & 0.0867 & 11.45\% \\
    全样本固定日内形状（事后口径，仅作下界） & 0.0964 & 12.72\% \\
    \midrule
    \multicolumn{3}{l}{\textit{日内对数偏差修正系数 }$\beta$\textit{ 的选取（对剩余时段的平均绝对误差）}} \\
    $\beta$ & 6:00 版 & 12:00 版 \\
    \midrule
    0 & 0.0480 & 0.0464 \\
    0.25 & 0.0465 & 0.0443 \\
    0.5 & \textbf{0.0458} & \textbf{0.0433} \\
    0.75 & 0.0460 & 0.0434 \\
    1 & 0.0469 & 0.0446 \\
    1.25 & 0.0485 & 0.0467 \\
    \bottomrule
  \end{tabular}
  \par\vspace{2pt}
  \begin{minipage}{.92\textwidth}\footnotesize
    注：``相对均价''以\textbf{评价区间（334 天）的均价} $0.7575$~元/kWh为基准，而非全年均价 $0.7662$~元/kWh。``全样本固定日内形状''使用了评价区间本身的均值，含未来信息，只是作为误差下界列出，不是可用策略。``近 $N$ 天滚动均值''系列表现不佳，原因见正文：把不同星期几混在一起平均，恰好抹掉了电价里唯一可利用的星期偏移。$\beta=0$ 即不做日内修正；$\beta=0.5$ 在三个更新时刻都最优或并列最优，$\beta\ge 1$ 时把已实现的偏差过度外推，反而变差。
  \end{minipage}
\end{table}

\begin{table}[htbp]
  \centering
  \caption{风险余量口径对照：以实际电价为权重的分位数 vs 普通分位数（正式区间按每 7 天抽样求均值）}
  \label{tab:p4-risk}
  \footnotesize
  \begin{tabular}{lrrrr}
    \toprule
    $\alpha$ & 价格加权/(kWh) & 普通分位数/(kWh) & 抬高量/(kWh) & 相对抬升 \\
    \midrule
    0.50 & -1.261 & -1.463 & +0.202 & --- \\
    0.60 & 12.210 & 11.700 & +0.510 & +4.36\% \\
    0.70 & 26.704 & 25.892 & +0.812 & +3.14\% \\
    0.80 & 44.026 & 43.047 & +0.979 & +2.27\% \\
    \bottomrule
  \end{tabular}
  \par\vspace{2pt}
  \begin{minipage}{.92\textwidth}\footnotesize
    注：$\alpha=0.50$ 处的余量为负，因为净负荷预测误差的中位数本身略小于零（日间预测略偏高），此时``相对抬升''无意义，故以 ``---'' 表示。抬高的绝对量约 $1$~kWh/时段，之所以不大，是因为同一时段的价格在日与日之间变动有限；但方向上它把余量推向高代价时段，且实现成本为零，故保留为主口径。
  \end{minipage}
\end{table}

\begin{table}[htbp]
  \centering
  \caption{净负荷预测误差 $\varepsilon$ 与实际电价的协方差检查（框架 3.2 节）}
  \label{tab:p4-cov}
  \footnotesize
  \begin{tabular}{lr}
    \toprule
    项目 & 数值 \\
    \midrule
    样本数 & 48\,096 \\
    合并口径相关系数 & $+0.0826$ \\
    逐时段相关系数均值 & $+0.206$ \\
    逐时段相关系数范围 & $+0.083$~$+0.401$ \\
    逐时段相关为正的比例 & 100.0\% \\
    \midrule
    电价最高 20\% 时段的 $\varepsilon$ 均值 & $+38.5$~kW \\
    电价最低 20\% 时段的 $\varepsilon$ 均值 & $-72.7$~kW \\
    两者之差 & $+111.2$~kW \\
    \bottomrule
  \end{tabular}
  \par\vspace{2pt}
  \begin{minipage}{.92\textwidth}\footnotesize
    注：合并口径把全部 (日期, 时段) 样本混在一起，被``时段之间''的电价差异（傍晚贵、凌晨便宜）主导；把时段固定住再看，相关系数由 $+0.083$ 升到 $+0.206$，且 100\% 的时段为正。这说明$\mathrm{Cov}(P,R)>0$ 是\textbf{同一时段内}的性质，故风险余量必须逐时段计算，不能对全年做一次汇总——这正是框架 3.4 节按 $t$ 单独求分位数的理由。
  \end{minipage}
\end{table}

\begin{table}[htbp]
  \centering
  \caption{波动电价下 4-3 的八种预报使用组合（各组合单独标定参数，评价区间与计费口径完全一致）}
  \label{tab:p4-combo}
  \footnotesize
  \begin{tabular}{lrrrrrr}
    \toprule
    使用的预报时刻 & 初始计划费/元 & 调整费/元 & 紧急费/元 & 合计费用/元 & 相对 $\varnothing$ & 期末储电量/kWh \\
    \midrule
    {6,12,18} & 13,266,140 & 675,425 & 255,250 & \textbf{14,196,815} & -5.46\% & 5,912.8 \\
    {6,18} & 13,375,475 & 636,131 & 222,257 & \textbf{14,233,863} & -5.21\% & 5,912.8 \\
    {6,12} & 13,457,237 & 495,800 & 488,102 & \textbf{14,441,138} & -3.83\% & 5,912.8 \\
    {12,18} & 13,547,270 & 420,986 & 491,980 & \textbf{14,460,237} & -3.71\% & 6,047.6 \\
    {18} & 13,693,592 & 365,481 & 405,543 & \textbf{14,464,617} & -3.68\% & 6,047.6 \\
    {12} & 13,704,099 & 278,654 & 649,240 & \textbf{14,631,993} & -2.56\% & 6,047.6 \\
    {6} & 13,916,529 & 222,144 & 576,868 & \textbf{14,715,541} & -2.01\% & 6,047.6 \\
    仅 0:00 & 14,218,788 & 0 & 798,139 & \textbf{15,016,926} & --- & 6,286.1 \\
    \bottomrule
  \end{tabular}
  \par\vspace{2pt}
  \begin{minipage}{.92\textwidth}\footnotesize
    注：与问题三的同一张表对照可以看出价格波动改变了什么——在固定电价下``用不用新预报''主要体现在调整费与紧急费的此消彼长；在波动电价下，更新预报同时改善了\textbf{电量}与\textbf{价格}两项判断（6:00 起既能看到当天更准的负载与光伏，也能用当天已实现的价格修正剩余时段的电价预测），因此增量收益的来源更宽。所有方案都使用 0:00 发布的预报，组合列出的时刻是额外使用的版本。
  \end{minipage}
\end{table}

\begin{table}[htbp]
  \centering
  \caption{两个分支在 334 个交付日上的检查结果（未通过天数均为零）}
  \label{tab:p4-check}
  \footnotesize
  \begin{tabular}{lrr}
    \toprule
    检查项 & result4-2 未通过天数 & result4-3 未通过天数 \\
    \midrule
    信息边界 & 0 & 0 \\
    计划冻结 & 0 & 0 \\
    物理可行性 & 0 & 0 \\
    状态累计 & 0 & 0 \\
    费用一致性 & 0 & 0 \\
    调整权限 & 0 & 0 \\
    \midrule
    \multicolumn{3}{l}{\textit{数值口径的最坏值}} \\
    \midrule
    逐段电量平衡最大残差/kWh & $0.00$ & $0.00$ \\
    计划费重算最大偏差/元 & $0.00$ & $0.00$ \\
    调整费重算最大偏差/元 & --- & $0.00$ \\
    紧急费重算最大偏差/元 & $0.00$ & $0.00$ \\
    跨日衔接最大偏差/kWh & $0.00$ & $0.00$ \\
    全局储电量范围/kWh & 1,200.0~10,800.0 & 1,200.0~10,800.0 \\
    \midrule
    储电下界余量/kWh & 0.000 & 0.000 \\
    储电上界余量/kWh & 0.000 & 0.000 \\
    \bottomrule
  \end{tabular}
  \par\vspace{2pt}
  \begin{minipage}{.92\textwidth}\footnotesize
    注：残差均为浮点误差量级（$<10^{-12}$~kWh），说明约束是被精确满足而非近似满足。储电量上下界余量恰为零，说明日末储备与容量约束在个别日期是\textbf{紧的}——这正是风险余量与储备线起作用的地方，不是缺陷。``计划冻结''在 result4-2 下的可核验内容是逐段电量平衡：执行层唯一的外部常规购电来源是 0:00 定下的计划量，日内不存在第二次常规购电。
  \end{minipage}
\end{table}
